\documentclass[10pt,a4paper]{amsart}
\usepackage[margin=19mm]{geometry}
\usepackage{amsmath,amssymb,mathtools}
\usepackage[scr=rsfs,cal=euler]{mathalfa}
\usepackage{xfrac}
\usepackage[unicode,hidelinks,bookmarks=false]{hyperref}
\newcommand{\ie}{i.e.\ }
\newcommand{\cf}{cf.\ }
\newcommand{\resp}{respectively\ }
\newcommand{\Subsection}{\subsection}
\providecommand{\nobreakdash}{\nobreak\hbox{-}}
\setlength{\emergencystretch}{2em}
\multlinegap0pt
\usepackage{bm}
\usepackage{bbm}
\usepackage{dsfont}
\usepackage{xspace}
\usepackage{cancel}
\usepackage{mathtools}
\usepackage{amsthm}
\makeatletter
\@ifundefined{theorem}{\newtheorem{theorem}{Theorem}}{}
\@ifundefined{corollary}{\newtheorem{corollary}[theorem]{Corollary}}{}
\makeatother
\newcommand{\R}{\mathbb{R}}
\newcommand{\U}{\mathbf{U}}
\newcommand{\e}{\varepsilon}
\title[Well-posedness and discrete-to-continuum convergence of the ]{Well-posedness and discrete-to-continuum convergence of the Kirchhoff Network model}
\author{Joaqu\'in Oyarz\'un, Marco Veneroni}
\date{Source: arXiv:2609.18351v1 (CC BY 4.0)}
\begin{document}
\maketitle

\noindent\emph{Source and licence.} Well-posedness and discrete-to-continuum convergence of the Kirchhoff Network model. Version arXiv:2609.18351v1 (CC BY 4.0), distributed under the Creative Commons Attribution 4.0 International licence, as recorded in the arXiv licence metadata for this exact version and shown on the abstract page https://arxiv.org/abs/2609.18351v1. Full rights notice: licenses/arxiv-2609.18351-CC-BY-4.0.txt.\par\smallskip

\noindent\textit{Editorial scope.} Counted unit: the Corollary whose source label is \texttt{cor:minimizers} in the article (source0.tex lines 1176--1196, source label \texttt{cor:minimizers}), delivered together with the article's own complete original proof of it (source0.tex lines 1200--1240, one \texttt{proof} environment of 41 lines). No mathematics is written, repaired or re-derived: statement and proof are the article's own text line for line. The proof is detached from the statement in the source: the article states the result at lines 1176--1196 and prints its proof at lines 1200--1240, and the proof environment's own header names the statement, so the association is the article's own. Required context retained uncounted: eq:FHN (lines 111--114); eq:zeromean (lines 121--124); eq:compatibility (lines 220--223); apriori-estimate (lines 233--241); th:main (lines 260--272); KNM-abstract (lines 431--441); MM-eps (lines 533--536); final-error-estimate (lines 623--629); Bidomain-abstract (lines 660--668); MM-eps-tt (lines 684--687); final-error-estimate-bidomain (lines 702--708); def:Hepsilon (lines 790--797); def:H (lines 810--817). Every definition, display and label of those lines is retained unchanged. Omitted from this package and not redistributed: 1--110, 115--120, 125--219, 224--232, 242--259, 273--430, 442--532, 537--622, 630--659, 669--683, 688--701, 709--789, 798--809, 818--1175, 1197--1199, 1241--1276. Nothing inside a retained excerpt is omitted or altered: every retained body is delivered exactly as the article's lines are written, so no omission receipt of any kind accompanies this package. Citations: the retained excerpts carry no citation command at all, so no bibliography entry of the article is transcribed and no reference is redistributed. Reference adaptation: none was needed and none was made; every reference inside the delivered unit resolves inside it, so no unresolved reference or citation is produced. Printed slips kept literal: no printed word, symbol, hypothesis or number of the counted statement or of its proof was altered. Layout and class: the article's own class is \texttt{article} with packages including babel, inputenc, amsthm, amsmath, amssymb, mathrsfs, textcomp, enumerate, setspace, latexsym, graphicx, mathtools, tabularray, geometry; the wrapper is the lane's pinned \texttt{amsart} editorial wrapper at 10pt on A4, which restates the theorem-like environments the retained text needs and adds the article's own macro definitions that the retained text needs (\texttt{\textbackslash R}, \texttt{\textbackslash U}, \texttt{\textbackslash e}), with extra wrapper packages (bm, bbm, dsfont, xspace, cancel, mathtools). The delivered numbering is the unit's own. Assets: no retained span contains a figure environment, a graphics-inclusion command, a PDF-page inclusion, a table or a TikZ picture; no image file is copied into this package, the package declares no asset dependency and no image file is redistributed with it. Licence: version arXiv:2609.18351v1 of this article is distributed under the Creative Commons Attribution 4.0 International licence, as recorded in the arXiv licence metadata for this exact version and shown on the abstract page https://arxiv.org/abs/2609.18351v1. A full rights notice is delivered at licenses/arxiv-2609.18351-CC-BY-4.0.txt. Characters: every retained excerpt is pure ASCII, so the delivered engine, XeLaTeX with its default Latin Modern text font, reports no missing character.
\begin{equation}
	\label{eq:FHN}
	I_{\text{ion}}(v, s) = F(v)+ \delta s, \qquad H(v,s) = \sigma v-\gamma s,
\end{equation}

\begin{equation}
	\label{eq:zeromean}
	\int_{\Omega} u_e(x, t) \, dx = 0, \quad \text{for a.e. } t \in (0,T).
\end{equation}

\begin{equation}
	\label{eq:compatibility}
	\bar u_0^\e \in \text{argmin}\left\{ a^\e(\bar u): u_i-u_e = v_0^\e\right\}. 
\end{equation}

	\begin{equation}
		\label{apriori-estimate}
		\begin{aligned}
			\sup_{t \in [0,T]} \left( b^\e(\bar{u}^\e(t)) + a^\e(\bar{u}^\e(t)) +\phi^\e(\bar{u}^\e(t)) +  |s^\e(t)|^2_\e  \right) 
			\hspace{4cm} \\
			+ \int_0^T \left( |\partial_t v^\e(t)|_\e^2 + |\partial_t s^\e(t)|_\e^2 \right) dt 
			\leq C \left( b^\e(\bar{u}_0^\e) + \phi^\e(\bar{u}_0^\e) + a^\e(\bar{u}_0^\e) +|s_0^\e|^2_\e \right),
		\end{aligned}
	\end{equation}

\begin{theorem}
	\label{th:main}
	Let $(\e_j)_{j\in \mathbb{N}}$ be a sequence of positive real numbers converging to $0$. Let $\mathbf{u}_0 := (u_{i,0}, u_{e,0}, s_0) \in H^1(\Omega) \times H_\circ^1(\Omega) \times L^2(\Omega)$. Suppose that $\mathbf{u}_0^{\varepsilon_j} := (u_{i,0}^{\varepsilon_j}, u_{e,0}^{\varepsilon_j}, s_0^{\varepsilon_j}) \in \mathbf{X}^{\varepsilon_j}$ is a sequence of initial data satisfying \eqref{eq:compatibility} such that, as $j \to \infty$, 
	\[
	\Pi_{\varepsilon_j}\mathbf{u}_0^{\varepsilon_j} \to \mathbf{u}_0\quad \text{strongly in } L^2(\Omega)^3\quad \text{and}\quad
	\limsup_{j\to \infty} a^{\varepsilon_j}(\bar u_0^{\varepsilon_j}) +\phi^{\varepsilon_j}(\bar u_0^{\varepsilon_j}) < \infty. 
	\]
	Then, there exists a subsequence $(\e_{j_k})$ and two positive-definite conductivity tensors $M_i, M_e \in L^\infty(\Omega;  \R^{3 \times 3}_{sym})$ such that the sequence of discrete solutions $\mathbf{u}^{\varepsilon_{j_k}}$ of the KNM$-\varepsilon_{j_k}$ system satisfies 
	\[
	\Pi_{\varepsilon_{j_k}}\mathbf{u}^{\varepsilon_{j_k}}(t) \to \mathbf{u}(t)\quad \text{strongly in } L^2(\Omega)^3
	\] 
	for every $t \in [0,T]$, where $\mathbf{u}$ is the unique solution of the macroscopic bidomain model (BD), complemented by \eqref{eq:FHN}--\eqref{eq:zeromean}, with conductivity tensors $M_i$ and $M_e$.
\end{theorem}

\begin{equation}
	\label{KNM-abstract}
	\left\{
	\begin{aligned}
		&\frac{d}{dt} \mathbf{b}^\e(\mathbf{u}^\e(t), \widehat{\mathbf{u}}) 
		+ \mathbf{a}^\e(\mathbf{u}^\e(t), \widehat{\mathbf{u}}) 
		+ 	\mathcal{F}^\e(\mathbf{u}^\e(t), \widehat{\mathbf{u}}) = \mathbf{g}^\e(\mathbf{u}^\e(t), \widehat{\mathbf{u}}), \\
		&\mathbf{b}^\e(\mathbf{u}^\e(0), \widehat{\mathbf{u}}) = \mathbf{b}^\e(\mathbf{u}_0^\e, \widehat{\mathbf{u}}),
	\end{aligned}
	\right.
\end{equation}

\begin{equation}
	\label{MM-eps}
	\Psi_{n}^{\e,\tau}(\mathbf{U}) = \frac{1}{2\tau} \mathbf{b}^\e\left(\mathbf{U} - \mathbf{U}_{n-1}^{\e,\tau}\right) + \frac{1}{2} \mathbf{a}^\e(\mathbf{U}) + \mathbf{\phi}^\e(\mathbf{U}) - \mathbf{g}^\e\left(\mathbf{U}_{n-1}^{\e,\tau}, \mathbf{U}\right).
\end{equation}

\begin{theorem}
	Let $\mathbf{u}^\e$ be the unique solution of the time-continuous problem \eqref{KNM-abstract} and let $\widetilde{\U}^{\e, \tau}$ be the piecewise linear interpolant of the solutions $\{\U^{\e,\tau}_n\}_{n=0}^N$ of the minimizing movement scheme \eqref{MM-eps}. Then, there exists $C > 0$ (independent of $\tau$ and $\e$) such that:
	\begin{equation}
		\label{final-error-estimate}
		\sup_{t \in [0, T]} \mathbf{b}^\e(\mathbf{u}^\e(t) - \widetilde{\mathbf{U}}^{\e, \tau}(t)) + \int_0^T \mathbf{a}^\e(\mathbf{u}^\e(t) - \overline{\mathbf{U}}^{\e, \tau}(t)) \, dt \le C \tau E^\e(\mathbf{u}_0^\e).
	\end{equation}
\end{theorem}

\begin{equation}
	\label{Bidomain-abstract}
	\left\{
	\begin{aligned}
		&\frac{d}{dt} \mathbf{b}(\mathbf{u}(t), \widehat{\mathbf{u}}) + \mathbf{a}(\mathbf{u}(t), \widehat{\mathbf{u}}) + \mathcal{F}(\mathbf{u}(t), \widehat{\mathbf{u}}) = \mathbf{g}(\mathbf{u}(t), \widehat{\mathbf{u}}), \\
		&\mathbf{b}(\mathbf{u}(0), \widehat{\mathbf{u}}) = \mathbf{b}(\mathbf{u}_0, \widehat{\mathbf{u}}),
	\end{aligned}
	\right.
\end{equation}

\begin{equation}
	\label{MM-eps-tt}
	\Psi_{n}^{\tau}(\mathbf{U}) := \frac{1}{2\tau}\, \mathbf{b}\left(\mathbf{U} - \mathbf{U}_{n-1}^\tau\right) + \frac{1}{2}\, \mathbf{a}(\mathbf{U}) + \phi(\mathbf{U}) - \mathbf{g}\left(\mathbf{U}_{n-1}^\tau, \mathbf{U}\right).
\end{equation}

\begin{theorem}
	Let $\mathbf{u}$ be the unique solution of \eqref{Bidomain-abstract} and let $\mathbf{U}^{\tau}$ be its piecewise linear interpolant. There exists a constant $C > 0$, independent of $\tau$, such that:
	\begin{equation}
		\label{final-error-estimate-bidomain}
		\sup_{t \in [0,T]} \mathbf{b}(\mathbf{u}(t) - \mathbf{U}^{\tau}(t)) + \int_0^T \mathbf{a}(\mathbf{u}(t) - \mathbf{U}^{\tau}(t)) \, dt \le C \tau (\mathbf{b}(\mathbf{u}_0) + \phi(\mathbf{u}_0) + \mathbf{a}(\mathbf{u}_0)).
	\end{equation}
\end{theorem}

\begin{equation}
	\label{def:Hepsilon}
	\mathcal{H}_\e(U_i,U_e):=
	\begin{cases}
		a_i^\e(u_i)+a_e^\e(u_e) &\text{if } U_i=\Pi_\e u_i,\ U_e=\Pi_\e u_e, \\
		+\infty &\text{otherwise,}
	\end{cases}
\end{equation}

	\begin{equation}
		\label{def:H}
		\mathcal{H}(u_i,u_e):=
		\begin{cases}
			\displaystyle \int_\Omega g_i(x,\nabla u_i(x)) +g_e(x,\nabla u_e(x))\ d x &\text{if } (u_i,u_e)\in H^1(\Omega)^2 \\
			+\infty &\text{otherwise.}
		\end{cases}
	\end{equation}

\begin{corollary}
\label{cor:minimizers}
Let $v^\e \in X^\e$ and $v\in L^2(\Omega)$ be such that
\(
\lim_{\e \to 0}|v^\e|_\e = \|v\|_{L^2(\Omega)}.
\) 
If $\bar u^\e = (u_i^\e,u_e^\e)$ is the unique solution of the minimum problem
\[
\min \left\{ \mathcal{H}^\e(\bar u): \bar u \in X^\e \times X_0^\e,\ u_i-u_e = v^\e \right\} 
\]
and $\limsup_{\e\to 0} \mathcal{H}^\e(\bar u^\e)<+\infty$, then
\[
\Pi_\e \bar u^\e \to \bar u\quad \text{strongly in }L^2(\Omega)^2\quad 
\text{and}\quad  \lim_{\e \to 0}\mathcal{H}^\e(\bar u^\e)=\mathcal{H}(\bar u),
\]
where $\bar u$ is the unique solution of the minimum problem
\[
\min \left\{ \mathcal{H}(\bar u): \bar u \in H^1(\Omega) \times H^1_\circ(\Omega),\ u_i-u_e = v \right\}.
\]
and $\mathcal{H}^\e,\mathcal{H}$ were defined in \eqref{def:Hepsilon} and \eqref{def:H}.
\end{corollary}

\begin{proof}[Proof of Theorem \ref{th:main}]
Let $\mathbf{u}^{\e_j} = (u_i^{\e_j},u_e^{\e_j},s^{\e_j})$ be a sequence of  solutions of the KNM$-\e_j$ problem, such that $\mathcal H_{\e_j}$ $\Gamma$-converges to $\mathcal H$. Let $M_i$ and $M_e$ be the conductivity tensors in $\mathcal H$ and let $\mathbf{u}=(u_i,u_e,s)$ be the solution of the corresponding BD model. Fix $t \in [0, T]$. Owing to estimates \eqref{apriori-estimate} and Corollary \ref{cor:minimizers}, we only have to show that $\Pi_{\varepsilon_j} v^{\varepsilon_j}(t)=\Pi_{\varepsilon_j}\left(u_i^{\varepsilon_j}(t) - u_e^{\varepsilon_j}(t)\right) \to v(t)$ and that $\Pi_{\varepsilon_j} s^{\varepsilon_j}(t) \to s(t)$, strongly in $L^2(\Omega)$. For any given time step size $\tau > 0$, we denote by $(V^{\varepsilon_j,\tau},S^{\varepsilon_j,\tau})$ and $(V^\tau,S^\tau)$ the time-dependent piecewise linear interpolants of the discrete solutions generated by the minimizing movements schemes \eqref{MM-eps} and \eqref{MM-eps-tt}, respectively. Dropping the variable $t$ for sake of notation, we have:
\begin{equation}
	\label{eq:error_decomp}
	\begin{aligned}
		& \|\Pi_{\varepsilon_j}(v^{\varepsilon_j},s^{\varepsilon_j}) - (v,s)\|_{L^2(\Omega)^2}\leq \, \| \Pi_{\varepsilon_j}(v^{\varepsilon_j},s^{\varepsilon_j}) - \Pi_{\varepsilon_j}(V^{\varepsilon_j,\tau},S^{\varepsilon_j,\tau})\|_{L^2(\Omega)^2} \\
		&\quad \quad + \|\Pi_{\varepsilon_j}(V^{\varepsilon_j,\tau},S^{\varepsilon_j,\tau}) - (V^{\tau},S^{\tau})\|_{L^2(\Omega)^2} 
		+ \|(V^{\tau},S^{\tau}) - (v,s)\|_{L^2(\Omega)^2}.
	\end{aligned}
\end{equation}
We begin by estimating the first and third terms on the right-hand side of \eqref{eq:error_decomp}. By hypothesis, the sequence of well-prepared microscopic initial data satisfies the uniform energy bound $\limsup_{j \to \infty} E^{\varepsilon_j}(\mathbf{u}_0^{\varepsilon_j}) = M < \infty$. Consequently, by \eqref{final-error-estimate}, we have
\begin{equation}
	\label{eq:c1tau}
	\begin{aligned}
		\| \Pi_{\varepsilon_j}(v^{\varepsilon_j},s^{\varepsilon_j}) - \Pi_{\varepsilon_j}(V^{\varepsilon_j,\tau},S^{\varepsilon_j,\tau})\|_{L^2(\Omega)^2}^2  &= | v^{\varepsilon_j} -V^{\varepsilon_j,\tau}|_{\varepsilon_j}^2 +| s^{\varepsilon_j}-S^{\varepsilon_j,\tau}|_{\varepsilon_j}^2, \\
		&\leq \max\left\{\frac{1}{C_m},1\right\}\mathbf{b}^{\varepsilon_j}(\mathbf{u}^{\varepsilon_j}-{\mathbf{U}}^{\varepsilon_j, \tau}),\\
		& \leq C_1^2 \tau,
	\end{aligned}
\end{equation}
where the positive constant $C_1$ depends exclusively on $M$, $C_m$, and $C$. Analogously, applying the corresponding continuous time-discretization error estimate \eqref{final-error-estimate-bidomain} we obtain
\begin{equation}
	\label{eq:c2tau}
	\|(V^{\tau},S^{\tau}) - (v,s)\|_{L^2(\Omega)^2}
	\leq C_2 \sqrt\tau.
\end{equation}
Next, we address the spatial limit $j \to \infty$ for the remaining central term of \eqref{eq:error_decomp} at the fixed time-step level $\tau > 0$. Let $n \in \{1, \dots, N\}$ denote the specific time interval containing the instant $t$. We carry out an inductive argument on the time step index $n$. For $n=0$, the strong convergence holds by the hypothesis on the initial data. Assume now that $\Pi_{\varepsilon_j} \mathbf{U}_{n-1}^{\varepsilon_j,\tau} \xrightarrow{s-s-w} \mathbf{U}_{n-1}^\tau$ as $j \to \infty$. Under this inductive step, Corollary 1 guarantees that the sequence of unique discrete minimizers $\mathbf{U}_n^{\varepsilon_j,\tau} = (U_{i,n}^{\varepsilon_j,\tau}, U_{e,n}^{\varepsilon_j,\tau}, S_n^{\varepsilon_j,\tau})$ satisfies
\begin{equation*}
	\Pi_{\varepsilon_j} \mathbf{U}_n^{\varepsilon_j,\tau} \xrightarrow{s-s-w} \mathbf{U}_n^\tau = (U_{i,n}^\tau, U_{e,n}^\tau, S_n^\tau) \quad \text{as } j \to \infty,
\end{equation*}
where $\mathbf{U}_n^\tau$ is the unique minimizer of the continuous functional $\Psi_n^\tau$. To upgrade the weak convergence $\Pi_{\varepsilon_j} S_n^{\varepsilon_j,\tau} \rightharpoonup S_n^\tau$ in $L^2(\Omega)$ to the strong topology, we exploit the convergence of $\mathbf{b}^{\varepsilon_j}(\mathbf{U}_n^{\varepsilon_j,\tau})$ guaranteed by Corollary 1 (iii), which implies that $\lim_{j \to \infty} \|\Pi_{\varepsilon_j} S^{\varepsilon_j,\tau}_n\|_{L^2(\Omega)} = \|S_n^\tau\|_{L^2(\Omega)}.$ 
Since the linear interpolant $\mathbf{U}^{\varepsilon_j, \tau}(t)$ is a convex combination of the discrete vectors $\mathbf{U}_{n-1}^{\varepsilon_j,\tau}$ and $\mathbf{U}_n^{\varepsilon_j,\tau}$ on each time subinterval, this strong convergence naturally extends pointwise in time to the instant $t$, yielding
\begin{equation}
	\label{eq:pescayogurt}
	\lim_{j \to \infty}\|\Pi_{\varepsilon_j}(V^{\varepsilon_j,\tau}(t),S^{\varepsilon_j,\tau}(t)) - (V^{\tau}(t),S^{\tau}(t))\|_{L^2(\Omega)^2} = 0.
\end{equation}
Passing to the limit superior as $j \to \infty$ on both sides of inequality \eqref{eq:error_decomp}, by \eqref{eq:c1tau}, \eqref{eq:c2tau}, and \eqref{eq:pescayogurt}, we obtain
\begin{equation*}
	\limsup_{j \to \infty}  \| \Pi_{\varepsilon_j} (v^{\varepsilon_j}(t),s^{\varepsilon_j}(t)) - (v(t),s(t))\|_{L^2(\Omega)^2}  \leq (C_1 + C_2) \sqrt{\tau}.
\end{equation*}
Taking the limit as the time step size $\tau \to 0^+$ concludes the proof.
\end{proof}

\end{document}
